\section{Hybrid AC/DC Distribution DER Planning Under Carbon and Peak \texorpdfstring{$N{-}1$}{N-1} Reliability Constraints}

This section defines a planning model for hybrid AC/DC distribution systems with distributed photovoltaic generation, distributed battery energy storage, and converter-coupled AC/DC power exchange. The model is intended for feeder-level planning rather than transmission expansion. Its distinguishing feature is the joint treatment of annual carbon control, peak-time reliability security, and extreme-event resilience: the normal-operation layer uses weighted representative operating states, the reliability layer uses a peak-snapshot \(N{-}1\) contingency set evaluated through an FMEA-style frequency-duration model, and the resilience layer uses extreme-state \(N{-}k\) scenarios with post-fault partitioning and island support.

The linear network model follows the sign conventions and AC/DC variable semantics of Section~\ref{sec:powerflow}. In particular, bus injections are taken as positive when they inject into the network, the AC feeder is represented by a P/Q LinDistFlow model with voltage-magnitude variables, and the DC subnetwork is written using conductance-based voltage differences.

\subsection{Study Object and Assumptions}

Consider a hybrid distribution system with AC bus set \(\mathcal{N}^{ac}\), DC bus set \(\mathcal{N}^{dc}\), AC branch set \(\mathcal{E}^{ac}\), DC branch set \(\mathcal{E}^{dc}\), and voltage-source converter set \(\mathcal{E}^{vsc}\). Let \(\mathcal{N}=\mathcal{N}^{ac}\cup\mathcal{N}^{dc}\). The planning horizon is one representative year.

Normal operation is represented by a weighted set of operating snapshots \(\Omega\), with weight \(\pi_\omega\ge 0\) and equivalent duration \(\Delta_\omega\) hours, satisfying
\begin{align}
\sum_{\omega\in\Omega} \pi_\omega = 1,
\qquad
\sum_{\omega\in\Omega} \pi_\omega\Delta_\omega = 8760.
\end{align}
Among these snapshots, let \(\omega^{pk}\in\Omega\) denote the system peak snapshot used for reliability planning.

Let \(\Xi\) denote a finite library of low-probability extreme states, such as storm, wildfire, flood, or heat-driven common-mode outage conditions. Each \(\xi\in\Xi\) has scenario weight \(\varpi_\xi\ge 0\), and for each extreme state let \(\mathcal{K}_\xi\) denote the associated \(N{-}k\) contingency family induced by that event.

Reliability is enforced through a single-contingency set \(\mathcal{C}\) covering the AC/DC components to be protected under an \(N{-}1\) criterion. Only the peak snapshot is analyzed under contingencies. For each failed component \(k\in\mathcal{C}\), let \(\lambda_k\) denote annual failure frequency and let \(\tau_k\) denote the equivalent outage duration used in the peak-time FMEA approximation.

The following assumptions define the intended model scope.

\begin{assumption}[Linear hybrid network model]
The planning model uses a linearized AC/DC power-flow approximation consistent with the variable conventions in Section~\ref{sec:powerflow}. The AC feeder is assumed radial or radially operated and is represented by a P/Q LinDistFlow approximation with branch active-power, branch reactive-power, and squared-voltage variables. DC active-power flow is represented by conductance-weighted DC-voltage differences, and each AC/DC converter is treated as a lossless controllable interlinking device.
\end{assumption}

\begin{assumption}[DER-assisted voltage regulation]
Reactive support may be supplied by PV inverters, BESS inverters, AC/DC converters, and upstream substations through linear box-type capability envelopes. Hence the AC voltage profile is regulated jointly by network reinforcement and controllable DER-side reactive power rather than by substations alone.
\end{assumption}

\begin{assumption}[Hierarchical reliability screening]
The reliability architecture has two security layers. The first evaluates only the peak snapshot \(\omega^{pk}\) under \(N{-}1\) contingencies. The second evaluates extreme states \(\xi\in\Xi\) under event-induced \(N{-}k\) contingencies, where post-fault partitioning and islanded operation are allowed.
\end{assumption}

\begin{assumption}[Tree-structured partition restoration]
Under extreme states, surviving substations and selected grid-forming BESS units may energize multiple supply partitions. Each energized partition must remain radial, so the post-fault topology is a forest rooted at activated substations or activated grid-forming BESS sources.
\end{assumption}

\begin{assumption}[FMEA weighting]
Reliability impact is measured through a frequency-duration-weighted peak load curtailment functional, consistent with the deterministic FMEA logic in Section~\ref{sec:distribution-reliability-fmea}.
\end{assumption}

\begin{assumption}[Snapshot storage approximation]
The planning model is not a chronological unit-commitment problem. Storage energy limits are therefore enforced by linear snapshot energy envelopes based on effective discharge/charge durations and a prescribed pre-contingency state-of-charge factor at the peak snapshot.
\end{assumption}

\subsection{Sets, Parameters, and Planning Variables}

\subsubsection{Network and device sets}

\begin{align}
\mathcal{R}^{ac} &\subseteq \mathcal{N}^{ac} && \text{AC root or substation buses}, \\
\mathcal{R}^{dc} &\subseteq \mathcal{N}^{dc} && \text{DC root or substation buses}, \\
\mathcal{N}^{pv} &\subseteq \mathcal{N} && \text{candidate PV buses}, \\
\mathcal{N}^{bess} &\subseteq \mathcal{N} && \text{candidate BESS buses}, \\
\mathcal{N}^{gfm} &\subseteq \mathcal{N}^{bess}\cap\mathcal{N}^{ac} && \text{candidate AC-side grid-forming BESS buses}, \\
\mathcal{N}^{crit} &\subseteq \mathcal{N} && \text{critical-load buses}, \\
\delta_i^{-},\delta_i^{+} &&& \text{incoming and outgoing incident branches of bus } i.
\end{align}

For each converter \(c\in\mathcal{E}^{vsc}\), let \(i(c)\in\mathcal{N}^{ac}\) be the associated AC bus and \(k(c)\in\mathcal{N}^{dc}\) be the associated DC bus.

For each AC branch \(e=(i,j)\in\mathcal{E}^{ac}\), the branch orientation is taken from the upstream bus \(i\) to the downstream bus \(j\), consistent with a radial feeder embedding used by LinDistFlow.

\subsubsection{Planning variables}

The first-stage decision vector is
\begin{align}
x := \Big(
&\overline P_i^{pv}\big)_{i\in\mathcal{N}^{pv}},
\big(\overline P_i^{bess},\overline E_i^{bess}\big)_{i\in\mathcal{N}^{bess}},
\big(\overline S_c^{vsc}\big)_{c\in\mathcal{E}^{vsc}}, \\
&\big(\overline P_r^{sub}\big)_{r\in\mathcal{R}^{ac}\cup\mathcal{R}^{dc}},
\zeta^{CO_2},\zeta^{rel},\zeta^{ext}
\Big),
\end{align}
where the components represent PV capacity, BESS power capacity, BESS energy capacity, converter capacity, substation import capability, carbon slack, peak-reliability slack, and extreme-resilience slack.

The annualized investment cost is written as
\begin{align}
C^{inv}(x) =
&\sum_{i\in\mathcal{N}^{pv}} \widehat K_i^{pv}\overline P_i^{pv}
+ \sum_{i\in\mathcal{N}^{bess}} \left(\widehat K_i^{bp}\overline P_i^{bess} + \widehat K_i^{be}\overline E_i^{bess}\right) \\
&+ \sum_{c\in\mathcal{E}^{vsc}} \widehat K_c^{vsc}\overline S_c^{vsc}
+ \sum_{r\in\mathcal{R}^{ac}\cup\mathcal{R}^{dc}} \widehat K_r^{sub}\overline P_r^{sub}.
\label{eq:10b-investment-cost}
\end{align}

\subsection{Three-Layer Planning Model}

The planning model contains three recourse layers:
\begin{enumerate}
  \item a normal-operation layer over \(\omega\in\Omega\), used for expected operating cost and annual carbon accounting;
  \item a peak-contingency layer over \(k\in\mathcal{C}\), used for \(N{-}1\) reliability enforcement by FMEA weighting;
  \item an extreme-state resilience layer over \((\xi,\kappa)\in\Xi\times\mathcal{K}_\xi\), used for \(N{-}k\) resilience enforcement under partitioned post-fault operation. In the current code base, this layer has both a legacy heuristic restoration simulator and a new strict multi-period MIP/LinDistFlow restoration skeleton for unexpected-fault replacement.
\end{enumerate}

The complete planning problem is
\begin{align}
\min_{x,\,y,\,z} \quad
& C^{inv}(x)
+ \sum_{\omega\in\Omega} \pi_\omega Q^{nor}(x,\omega)
+ \rho_{CO_2}\zeta^{CO_2}
+ \rho_{rel}\zeta^{rel}
+ \rho_{ext}\zeta^{ext}
\label{eq:10b-master-obj}\\
\text{s.t.}\quad
& x\in\mathcal X,
\label{eq:10b-master-xset}\\
& \sum_{\omega\in\Omega} \pi_\omega G^{CO_2}(x,\omega) \le \overline E^{ann}_{CO_2} + \zeta^{CO_2},
\label{eq:10b-master-carbon}\\
& \sum_{k\in\mathcal{C}} \lambda_k\tau_k H^{pk}_k(x) \le \overline E^{pk}_{rel} + \zeta^{rel},
\label{eq:10b-master-reliability}\\
& \sum_{\xi\in\Xi} \varpi_\xi \sum_{\kappa\in\mathcal{K}_\xi} \lambda_{\xi,\kappa}\tau_{\xi,\kappa} H^{ext}_{\xi,\kappa}(x)
\le \overline E^{ext}_{rel} + \zeta^{ext},
\label{eq:10b-master-extreme}\\
& \zeta^{CO_2} \ge 0,
\qquad
\zeta^{rel} \ge 0,
\qquad
\zeta^{ext} \ge 0.
\label{eq:10b-master-slacks}
\end{align}

Here \(Q^{nor}(x,\omega)\) is normal-state operating cost, \(G^{CO_2}(x,\omega)\) is snapshot carbon contribution, \(H_k^{pk}(x)\) is peak-snapshot unsupplied energy under contingency \(k\), and \(H^{ext}_{\xi,\kappa}(x)\) is unsupplied energy under extreme state \(\xi\) and higher-order contingency \(\kappa\). The scalar \(\overline E^{ann}_{CO_2}\) is the annual carbon budget, \(\overline E^{pk}_{rel}\) is an admissible FMEA-weighted peak EENS threshold, and \(\overline E^{ext}_{rel}\) is the resilience budget for extreme-state \(N{-}k\) events.

\subsection{Normal-Operation Recourse Model}

\subsubsection{Operating variables}

For each operating snapshot \(\omega\in\Omega\), the normal-state recourse variables are
\begin{align}
y_\omega := \Big(
&f_{e,\omega}^{ac},q_{e,\omega}^{ac},u_{i,\omega}^{ac},
f_{\ell,\omega}^{dc},v_{k,\omega}^{dc},
p_{c,\omega}^{vsc},q_{c,\omega}^{vsc}, \\
&p_{i,\omega}^{pv},q_{i,\omega}^{pv},
p_{i,\omega}^{bess,ch},p_{i,\omega}^{bess,dis},q_{i,\omega}^{bess},
p_{r,\omega}^{sub,+},p_{r,\omega}^{sub,-},q_{r,\omega}^{sub}
\Big),
\end{align}
where \(p_{c,\omega}^{vsc}\in\mathbb R\) is the signed converter transfer. A positive value denotes active-power transfer from the DC side to the AC side, while a negative value denotes transfer from the AC side to the DC side.

\subsubsection{Operating objective}

For fixed \((x,\omega)\), define
\begin{align}
Q^{nor}(x,\omega) := 
\min_{y_\omega} \quad
& \sum_{r\in\mathcal{R}^{ac}\cup\mathcal{R}^{dc}}
\left(c_{r,\omega}^{buy} p_{r,\omega}^{sub,+} - c_{r,\omega}^{sell} p_{r,\omega}^{sub,-}\right) \\
&+ \sum_{i\in\mathcal{N}^{bess}} c_i^{bess}
\left(p_{i,\omega}^{bess,ch}+p_{i,\omega}^{bess,dis}\right)
\label{eq:10b-normal-obj}
\end{align}
subject to the following constraints.

\paragraph{PV and storage capability.}
For all \(i\in\mathcal{N}^{pv}\cap\mathcal{N}^{bess}\) and \(\omega\in\Omega\),
\begin{align}
0 \le p_{i,\omega}^{pv} &\le \alpha_{i,\omega}^{pv}\overline P_i^{pv},
\label{eq:10b-pv-limit}\\
-\overline\gamma_i^{pv}\alpha_{i,\omega}^{pv}\overline P_i^{pv}
\le q_{i,\omega}^{pv} \le
\overline\gamma_i^{pv}\alpha_{i,\omega}^{pv}\overline P_i^{pv},
\label{eq:10b-pv-reactive-limit}\\
0 \le p_{i,\omega}^{bess,ch} &\le \overline P_i^{bess},
\label{eq:10b-bess-charge}\\
0 \le p_{i,\omega}^{bess,dis} &\le \overline P_i^{bess},
\label{eq:10b-bess-discharge}\\
-\overline\gamma_i^{bess}\overline P_i^{bess}
\le q_{i,\omega}^{bess} \le
\overline\gamma_i^{bess}\overline P_i^{bess},
\label{eq:10b-bess-reactive-limit}\\
p_{i,\omega}^{bess,ch}\Delta_\omega &\le (1-\underline s_i)\overline E_i^{bess},
\label{eq:10b-bess-energy-charge}\\
p_{i,\omega}^{bess,dis}\Delta_\omega &\le \overline s_i\overline E_i^{bess}.
\label{eq:10b-bess-energy-discharge}
\end{align}
If a candidate bus hosts only PV or only BESS, the irrelevant variables are omitted.

\paragraph{Substation capacity.}
For each root bus \(r\),
\begin{align}
0 \le p_{r,\omega}^{sub,+} \le \overline P_r^{sub},
\qquad
0 \le p_{r,\omega}^{sub,-} \le \overline P_r^{sub}.
\label{eq:10b-substation-bounds}
\end{align}
For each AC root bus \(r\in\mathcal{R}^{ac}\),
\begin{align}
-\overline Q_r^{sub} \le q_{r,\omega}^{sub} \le \overline Q_r^{sub}.
\label{eq:10b-substation-reactive-bounds}
\end{align}

\paragraph{Lossless converter capability.}
For each converter \(c\in\mathcal{E}^{vsc}\),
\begin{align}
 -\overline S_c^{vsc} \le p_{c,\omega}^{vsc} \le \overline S_c^{vsc}.
\label{eq:10b-vsc-bounds}
\end{align}
On the AC side, the converter may also provide reactive support within a linear capability box:
\begin{align}
-\overline\gamma_c^{vsc}\overline S_c^{vsc}
\le q_{c,\omega}^{vsc} \le
\overline\gamma_c^{vsc}\overline S_c^{vsc}.
\label{eq:10b-vsc-reactive-bounds}
\end{align}
Under the lossless approximation, the converter contributes equal and opposite injections on the two subnetworks. Define
\begin{align}
p_{c,\omega}^{vsc,ac} &= p_{c,\omega}^{vsc},
\label{eq:10b-vsc-ac-inj}\\
p_{c,\omega}^{vsc,dc} &= -p_{c,\omega}^{vsc}.
\label{eq:10b-vsc-dc-inj}
\end{align}
The converter reactive injection appears only on the AC side and is denoted by
\begin{align}
q_{c,\omega}^{vsc,ac} = q_{c,\omega}^{vsc}.
\label{eq:10b-vsc-ac-reactive-inj}
\end{align}
Hence the converter only transfers active power between the AC and DC subnetworks, while AC-side reactive support is supplied locally by its converter station.

\paragraph{Linear AC P/Q LinDistFlow constraints.}
For each AC branch \(e=(i,j)\in\mathcal{E}^{ac}\),
\begin{align}
u_{j,\omega}^{ac} &= u_{i,\omega}^{ac} - 2\left(r_e^{ac} f_{e,\omega}^{ac} + x_e^{ac} q_{e,\omega}^{ac}\right),
\label{eq:10b-ac-flow}\\
-\overline F_e^{ac} \le f_{e,\omega}^{ac} &\le \overline F_e^{ac},
\label{eq:10b-ac-flow-bound}\\
-\overline Q_e^{ac} \le q_{e,\omega}^{ac} &\le \overline Q_e^{ac},
\label{eq:10b-ac-reactive-flow-bound}\\
\underline u_i^{ac} \le u_{i,\omega}^{ac} &\le \overline u_i^{ac}.
\label{eq:10b-ac-voltage-bound}
\end{align}
For each AC root bus \(r\in\mathcal{R}^{ac}\), the feeder head is fixed at
\begin{align}
u_{r,\omega}^{ac} = u_r^{ref}.
\label{eq:10b-ac-root-voltage}
\end{align}
Here \(u_{i,\omega}^{ac}\) denotes the squared voltage magnitude at AC bus \(i\). The P/Q LinDistFlow approximation retains both active- and reactive-power contributions to voltage drop, while neglecting the higher-order current and loss terms of the full nonlinear DistFlow equations.

\paragraph{Linear DC branch flow.}
For each DC branch \(\ell=(i,j)\in\mathcal{E}^{dc}\),
\begin{align}
f_{\ell,\omega}^{dc} &= g_\ell^{dc}(v_{i,\omega}^{dc}-v_{j,\omega}^{dc}),
\label{eq:10b-dc-flow}\\
-\overline F_\ell^{dc} \le f_{\ell,\omega}^{dc} &\le \overline F_\ell^{dc},
\label{eq:10b-dc-flow-bound}\\
\underline v_i^{dc} \le v_{i,\omega}^{dc} &\le \overline v_i^{dc}.
\label{eq:10b-dc-voltage-bound}
\end{align}

\paragraph{AC nodal balance.}
For each AC bus \(i\in\mathcal{N}^{ac}\),
\begin{align}
\sum_{e\in\delta_i^{-}\cap\mathcal{E}^{ac}} f_{e,\omega}^{ac}
- \sum_{e\in\delta_i^{+}\cap\mathcal{E}^{ac}} f_{e,\omega}^{ac}
+ p_{i,\omega}^{inj,ac} = 0,
\label{eq:10b-ac-balance}
\end{align}
where the net AC injection is
\begin{align}
p_{i,\omega}^{inj,ac} :=
&\mathbf{1}_{\{i\in\mathcal{N}^{pv}\}}p_{i,\omega}^{pv}
+ \mathbf{1}_{\{i\in\mathcal{N}^{bess}\}}\left(p_{i,\omega}^{bess,dis}-p_{i,\omega}^{bess,ch}\right) \\
&+ \mathbf{1}_{\{i\in\mathcal{R}^{ac}\}}\left(p_{i,\omega}^{sub,+}-p_{i,\omega}^{sub,-}\right)
+ \sum_{c:i(c)=i} p_{c,\omega}^{vsc,ac}
- P_{i,\omega}^{load,ac}.
\label{eq:10b-ac-net-injection}
\end{align}
For each AC bus \(i\in\mathcal{N}^{ac}\), reactive balance is enforced by
\begin{align}
\sum_{e\in\delta_i^{-}\cap\mathcal{E}^{ac}} q_{e,\omega}^{ac}
- \sum_{e\in\delta_i^{+}\cap\mathcal{E}^{ac}} q_{e,\omega}^{ac}
+ q_{i,\omega}^{inj,ac} = 0,
\label{eq:10b-ac-reactive-balance}
\end{align}
where
\begin{align}
q_{i,\omega}^{inj,ac} :=
\mathbf{1}_{\{i\in\mathcal{N}^{pv}\}} q_{i,\omega}^{pv}
+ \mathbf{1}_{\{i\in\mathcal{N}^{bess}\}} q_{i,\omega}^{bess}
+ \mathbf{1}_{\{i\in\mathcal{R}^{ac}\}} q_{i,\omega}^{sub}
+ \sum_{c:i(c)=i} q_{c,\omega}^{vsc,ac}
- Q_{i,\omega}^{load,ac}.
\label{eq:10b-ac-reactive-net-injection}
\end{align}
This equation explicitly captures DER-assisted voltage regulation through PV, BESS, VSC, and substation reactive-power decisions.

\paragraph{DC nodal balance.}
For each DC bus \(k\in\mathcal{N}^{dc}\),
\begin{align}
\sum_{\ell\in\delta_k^{-}\cap\mathcal{E}^{dc}} f_{\ell,\omega}^{dc}
- \sum_{\ell\in\delta_k^{+}\cap\mathcal{E}^{dc}} f_{\ell,\omega}^{dc}
+ p_{k,\omega}^{inj,dc} = 0,
\label{eq:10b-dc-balance}
\end{align}
where
\begin{align}
p_{k,\omega}^{inj,dc} :=
&\mathbf{1}_{\{k\in\mathcal{N}^{pv}\}}p_{k,\omega}^{pv}
+ \mathbf{1}_{\{k\in\mathcal{N}^{bess}\}}\left(p_{k,\omega}^{bess,dis}-p_{k,\omega}^{bess,ch}\right) \\
&+ \mathbf{1}_{\{k\in\mathcal{R}^{dc}\}}\left(p_{k,\omega}^{sub,+}-p_{k,\omega}^{sub,-}\right)
+ \sum_{c:k(c)=k} p_{c,\omega}^{vsc,dc}
- P_{k,\omega}^{load,dc}.
\label{eq:10b-dc-net-injection}
\end{align}

\subsubsection{Carbon accounting}

The annual carbon functional is defined from normal-state substation imports:
\begin{align}
G^{CO_2}(x,\omega) :=
\Delta_\omega
\sum_{r\in\mathcal{R}^{ac}\cup\mathcal{R}^{dc}}
\kappa_{r,\omega} p_{r,\omega}^{sub,+},
\label{eq:10b-carbon-functional}
\end{align}
where \(\kappa_{r,\omega}\) is the emission factor associated with upstream supply at root bus \(r\) during snapshot \(\omega\). Export is not credited unless a policy-adjusted negative emission factor is explicitly adopted.

\subsection{Peak-Time \texorpdfstring{$N{-}1$}{N-1} Reliability Model}

\subsubsection{Contingency variables}

For each contingency \(k\in\mathcal{C}\), the peak-time recourse vector is
\begin{align}
z_k := \Big(
&f_{e,k}^{ac},q_{e,k}^{ac},u_{i,k}^{ac},
f_{\ell,k}^{dc},v_{m,k}^{dc},
p_{c,k}^{vsc},q_{c,k}^{vsc}, \\
&p_{i,k}^{pv},q_{i,k}^{pv},
p_{i,k}^{bess,ch},p_{i,k}^{bess,dis},q_{i,k}^{bess},
p_{r,k}^{sub,+},p_{r,k}^{sub,-},q_{r,k}^{sub},
p_{n,k}^{ls}
\Big),
\end{align}
where \(p_{n,k}^{ls}\ge 0\) is curtailed load at peak bus \(n\) under contingency \(k\).

Let \(\zeta_{e,k}^{ac}\), \(\zeta_{\ell,k}^{dc}\), and \(\zeta_{c,k}^{vsc}\) be binary availability parameters equal to zero if the corresponding asset fails in contingency \(k\), and one otherwise.

\subsubsection{Peak-contingency constraints}

The peak-time post-contingency model reuses the normal-state equations with the peak data \((P^{load}_{\omega^{pk}},Q^{load}_{\omega^{pk}},\alpha^{pv}_{\omega^{pk}})\), but modifies device bounds by contingency availability.

\paragraph{Contingency network limits.}
For each AC branch \(e\in\mathcal{E}^{ac}\), DC branch \(\ell\in\mathcal{E}^{dc}\), and converter \(c\in\mathcal{E}^{vsc}\),
\begin{align}
 -M_e^{ac}(1-\zeta_{e,k}^{ac}) \le u_{j,k}^{ac} - u_{i,k}^{ac} + 2\left(r_e^{ac} f_{e,k}^{ac} + x_e^{ac} q_{e,k}^{ac}\right) \le M_e^{ac}(1-\zeta_{e,k}^{ac}),
\label{eq:10b-cont-ac-drop}\\
-\zeta_{e,k}^{ac}\overline F_e^{ac} \le f_{e,k}^{ac} \le \zeta_{e,k}^{ac}\overline F_e^{ac},
\label{eq:10b-cont-ac-flow}\\
-\zeta_{e,k}^{ac}\overline Q_e^{ac} \le q_{e,k}^{ac} \le \zeta_{e,k}^{ac}\overline Q_e^{ac},
\label{eq:10b-cont-ac-reactive-flow}\\
\underline u_i^{ac} \le u_{i,k}^{ac} \le \overline u_i^{ac},
\label{eq:10b-cont-ac-voltage}\\
 -M_\ell^{dc}(1-\zeta_{\ell,k}^{dc}) \le f_{\ell,k}^{dc} - g_\ell^{dc}(v_{i,k}^{dc}-v_{j,k}^{dc}) \le M_\ell^{dc}(1-\zeta_{\ell,k}^{dc}),
\label{eq:10b-cont-dc-drop}\\
-\zeta_{\ell,k}^{dc}\overline F_\ell^{dc} \le f_{\ell,k}^{dc} \le \zeta_{\ell,k}^{dc}\overline F_\ell^{dc},
\label{eq:10b-cont-dc-flow}\\
 -\zeta_{c,k}^{vsc}\overline S_c^{vsc} \le p_{c,k}^{vsc} \le \zeta_{c,k}^{vsc}\overline S_c^{vsc},
\label{eq:10b-cont-vsc-bound}\\
-\zeta_{c,k}^{vsc}\overline\gamma_c^{vsc}\overline S_c^{vsc}
\le q_{c,k}^{vsc} \le
\zeta_{c,k}^{vsc}\overline\gamma_c^{vsc}\overline S_c^{vsc}.
\label{eq:10b-cont-vsc-reactive-bound}
\end{align}
For each AC root bus \(r\in\mathcal{R}^{ac}\),
\begin{align}
u_{r,k}^{ac} = u_r^{ref},
\label{eq:10b-cont-ac-root-voltage}\\
-\overline Q_r^{sub} \le q_{r,k}^{sub} \le \overline Q_r^{sub}.
\label{eq:10b-cont-substation-reactive-bounds}
\end{align}
The big-\(M\) terms decouple post-contingency voltage-drop equations from outaged AC or DC branches. This prevents a failed line from imposing an artificial algebraic tie between the voltages of buses that should no longer be connected by that asset.

\paragraph{Peak PV and BESS envelopes.}
For all candidate DER buses,
\begin{align}
0 \le p_{i,k}^{pv} &\le \alpha_{i,\omega^{pk}}^{pv}\overline P_i^{pv},
\label{eq:10b-cont-pv}\\
-\overline\gamma_i^{pv}\alpha_{i,\omega^{pk}}^{pv}\overline P_i^{pv}
\le q_{i,k}^{pv} \le
\overline\gamma_i^{pv}\alpha_{i,\omega^{pk}}^{pv}\overline P_i^{pv},
\label{eq:10b-cont-pv-reactive}\\
0 \le p_{i,k}^{bess,ch} &\le \overline P_i^{bess},
\qquad
0 \le p_{i,k}^{bess,dis} \le \overline P_i^{bess},
\label{eq:10b-cont-bess-power}\\
-\overline\gamma_i^{bess}\overline P_i^{bess}
\le q_{i,k}^{bess} \le
\overline\gamma_i^{bess}\overline P_i^{bess},
\label{eq:10b-cont-bess-reactive}\\
p_{i,k}^{bess,dis}\tau_k &\le s_i^{pk}\overline E_i^{bess},
\qquad
p_{i,k}^{bess,ch}\tau_k \le (1-s_i^{pk})\overline E_i^{bess},
\label{eq:10b-cont-bess-energy}
\end{align}
where \(s_i^{pk}\in[\underline s_i,\overline s_i]\) is a prescribed pre-contingency state-of-charge factor at the peak snapshot.

\paragraph{Peak nodal balance with curtailment.}
For each AC bus \(i\in\mathcal{N}^{ac}\), active and reactive balances are enforced by
\begin{align}
\sum_{e\in\delta_i^{-}\cap\mathcal{E}^{ac}} f_{e,k}^{ac}
- \sum_{e\in\delta_i^{+}\cap\mathcal{E}^{ac}} f_{e,k}^{ac}
+ p_{i,k}^{inj,pk}
+ p_{i,k}^{ls}
= 0,
\label{eq:10b-cont-ac-balance}\\
\sum_{e\in\delta_i^{-}\cap\mathcal{E}^{ac}} q_{e,k}^{ac}
- \sum_{e\in\delta_i^{+}\cap\mathcal{E}^{ac}} q_{e,k}^{ac}
+ q_{i,k}^{inj,pk}
+ \eta_i^{Q,load} p_{i,k}^{ls}
= 0,
\label{eq:10b-cont-ac-reactive-balance}
\end{align}
where the peak active and reactive injections are
\begin{align}
p_{i,k}^{inj,pk} :=
&\mathbf{1}_{\{i\in\mathcal{N}^{pv}\}}p_{i,k}^{pv}
+ \mathbf{1}_{\{i\in\mathcal{N}^{bess}\}}\left(p_{i,k}^{bess,dis}-p_{i,k}^{bess,ch}\right) \\
&+ \mathbf{1}_{\{i\in\mathcal{R}^{ac}\}}\left(p_{i,k}^{sub,+}-p_{i,k}^{sub,-}\right)
+ \sum_{c:i(c)=i} p_{c,k}^{vsc,ac}
- P_{i,\omega^{pk}}^{load,ac},
\label{eq:10b-cont-ac-net-injection}\\
q_{i,k}^{inj,pk} :=
&\mathbf{1}_{\{i\in\mathcal{N}^{pv}\}} q_{i,k}^{pv}
+ \mathbf{1}_{\{i\in\mathcal{N}^{bess}\}} q_{i,k}^{bess}
+ \mathbf{1}_{\{i\in\mathcal{R}^{ac}\}} q_{i,k}^{sub}
+ \sum_{c:i(c)=i} q_{c,k}^{vsc}
- Q_{i,\omega^{pk}}^{load,ac}.
\label{eq:10b-cont-ac-reactive-net-injection}
\end{align}
For each DC bus \(m\in\mathcal{N}^{dc}\), active balance remains
\begin{align}
\sum_{\ell\in\delta_m^{-}\cap\mathcal{E}^{dc}} f_{\ell,k}^{dc}
- \sum_{\ell\in\delta_m^{+}\cap\mathcal{E}^{dc}} f_{\ell,k}^{dc}
+ p_{m,k}^{inj,pk}
+ p_{m,k}^{ls}
= 0,
\label{eq:10b-cont-dc-balance}
\end{align}
with
\begin{align}
p_{m,k}^{inj,pk} :=
&\mathbf{1}_{\{m\in\mathcal{N}^{pv}\}}p_{m,k}^{pv}
+ \mathbf{1}_{\{m\in\mathcal{N}^{bess}\}}\left(p_{m,k}^{bess,dis}-p_{m,k}^{bess,ch}\right) \\
&+ \mathbf{1}_{\{m\in\mathcal{R}^{dc}\}}\left(p_{m,k}^{sub,+}-p_{m,k}^{sub,-}\right)
+ \sum_{c:k(c)=m} p_{c,k}^{vsc,dc}
- P_{m,\omega^{pk}}^{load,dc}.
\label{eq:10b-cont-dc-net-injection}
\end{align}
The coefficient \(\eta_i^{Q,load}\) is the fixed reactive-to-active load ratio used to translate active load shedding into proportional reactive-demand reduction at AC bus \(i\).

\paragraph{Critical-load \texorpdfstring{$N{-}1$}{N-1} security.}
For each critical-load bus \(n\in\mathcal{N}^{crit}\) and each contingency \(k\in\mathcal{C}\),
\begin{align}
p_{n,k}^{ls} = 0.
\label{eq:10b-critical-n-minus-1}
\end{align}
This is the strict \(N{-}1\) condition for the critical load set. If full-load \(N{-}1\) service is required, \eqref{eq:10b-critical-n-minus-1} may be imposed for all buses.

\subsubsection{FMEA-weighted reliability metric}

The peak-contingency unsupplied-energy functional is
\begin{align}
H_k^{pk}(x) := \sum_{n\in\mathcal{N}} p_{n,k}^{ls}.
\label{eq:10b-peak-shed-functional}
\end{align}
The reliability constraint \eqref{eq:10b-master-reliability} therefore becomes
\begin{align}
\sum_{k\in\mathcal{C}} \lambda_k\tau_k
\sum_{n\in\mathcal{N}} p_{n,k}^{ls}
\le \overline E^{pk}_{rel} + \zeta^{rel}.
\label{eq:10b-fmea-peak-eens}
\end{align}
Equation \eqref{eq:10b-fmea-peak-eens} is a peak-only FMEA proxy for expected annual unserved energy. It is more informative than a pure deterministic no-shed rule because it retains failure-frequency and repair-duration information, yet it is far cheaper than full chronological contingency simulation.

\subsection{Extreme-State \texorpdfstring{$N{-}k$}{N-k} Resilience Model}

The third recourse layer addresses low-probability, high-impact states \(\xi\in\Xi\) in which weather, wildfire, flood, or cascading damage can induce simultaneous outages. For each extreme state \(\xi\), let \(\mathcal{K}_\xi\) denote the associated \(N{-}k\) contingency family. Unlike the peak \(N{-}1\) layer, the extreme-state layer permits network reconfiguration, multiple supply partitions, and grid-forming operation of selected BESS units.

\subsubsection{Extreme-state variables}

For each \((\xi,\kappa)\in\Xi\times\mathcal{K}_\xi\), define the resilience recourse vector
\begin{align}
w_{\xi,\kappa} := \Big(
&\beta_{a,\xi,\kappa},F_{a,\xi,\kappa},y_{n,\xi,\kappa}^{en},z_{s,\xi,\kappa}^{src},z_{i,\xi,\kappa}^{gfm}, \\
&f_{e,\xi,\kappa}^{ac},q_{e,\xi,\kappa}^{ac},u_{i,\xi,\kappa}^{ac},
f_{\ell,\xi,\kappa}^{dc},v_{m,\xi,\kappa}^{dc}, \\
&p_{c,\xi,\kappa}^{vsc},q_{c,\xi,\kappa}^{vsc},
p_{i,\xi,\kappa}^{pv},q_{i,\xi,\kappa}^{pv},
p_{i,\xi,\kappa}^{bess,ch},p_{i,\xi,\kappa}^{bess,dis},q_{i,\xi,\kappa}^{bess}, \\
&p_{r,\xi,\kappa}^{sub,+},p_{r,\xi,\kappa}^{sub,-},q_{r,\xi,\kappa}^{sub},
p_{n,\xi,\kappa}^{ls}
\Big),
\end{align}
where \(\beta_{a,\xi,\kappa}\in\{0,1\}\) is the energized status of hybrid network asset \(a\in\mathcal{E}^{ac}\cup\mathcal{E}^{dc}\cup\mathcal{E}^{vsc}\), \(y_{n,\xi,\kappa}^{en}\in\{0,1\}\) indicates whether bus \(n\) is energized, \(z_{s,\xi,\kappa}^{src}\in\{0,1\}\) activates a source partition rooted at substation or grid-forming BESS source \(s\), and \(z_{i,\xi,\kappa}^{gfm}\in\{0,1\}\) activates BESS \(i\in\mathcal{N}^{gfm}\) in grid-forming mode.

\subsubsection{Partition topology and source activation}

Let \(\mathcal{S}^{src} := \mathcal{R}^{ac}\cup\mathcal{R}^{dc}\cup\mathcal{N}^{gfm}\) be the candidate source set. Introduce a super-root node \(0\) and virtual source arcs \(\mathcal{E}^{sup}=\{(0,s): s\in\mathcal{S}^{src}\}\). Then tree-structured partitioning is imposed on the augmented graph by
\begin{align}
\beta_{a,\xi,\kappa} &\le \zeta_{a,\xi,\kappa},
\qquad
\forall a\in\mathcal{E}^{ac}\cup\mathcal{E}^{dc}\cup\mathcal{E}^{vsc},
\label{eq:10b-ext-availability}\\
\sum_{a}\beta_{a,\xi,\kappa} + \sum_{s\in\mathcal{S}^{src}} z_{s,\xi,\kappa}^{src}
&= \sum_{n\in\mathcal{N}} y_{n,\xi,\kappa}^{en},
\label{eq:10b-ext-tree-cardinality}\\
0 \le F_{a,\xi,\kappa} &\le M_a^{F}\beta_{a,\xi,\kappa},
\qquad
\forall a\in\mathcal{E}^{ac}\cup\mathcal{E}^{dc}\cup\mathcal{E}^{vsc},
\label{eq:10b-ext-flow-capacity}\\
0 \le F_{0s,\xi,\kappa} &\le M_s^{F} z_{s,\xi,\kappa}^{src},
\qquad
\forall s\in\mathcal{S}^{src},
\label{eq:10b-ext-source-arc}\\
\sum_{a\in\delta_n^{-}\cap(\mathcal{E}\cup\mathcal{E}^{sup})} F_{a,\xi,\kappa}
- \sum_{a\in\delta_n^{+}\cap(\mathcal{E}\cup\mathcal{E}^{sup})} F_{a,\xi,\kappa}
&= y_{n,\xi,\kappa}^{en},
\qquad
\forall n\in\mathcal{N},
\label{eq:10b-ext-connectivity}\\
\sum_{s\in\mathcal{S}^{src}} F_{0s,\xi,\kappa}
&= \sum_{n\in\mathcal{N}} y_{n,\xi,\kappa}^{en}.
\label{eq:10b-ext-super-root-balance}
\end{align}
Equations \eqref{eq:10b-ext-tree-cardinality}--\eqref{eq:10b-ext-super-root-balance} guarantee that the energized post-fault topology is a forest of radial partitions, each connected to exactly one activated source.

For grid-forming BESS activation,
\begin{align}
z_{i,\xi,\kappa}^{src} &\le z_{i,\xi,\kappa}^{gfm} \le y_{i,\xi,\kappa}^{en},
\qquad \forall i\in\mathcal{N}^{gfm},
\label{eq:10b-ext-gfm-source-link}\\
0 \le p_{i,\xi,\kappa}^{bess,dis} &\le \overline P_i^{gfm} z_{i,\xi,\kappa}^{gfm},
\qquad \forall i\in\mathcal{N}^{gfm},
\label{eq:10b-ext-gfm-p}\\
-\overline Q_i^{gfm} z_{i,\xi,\kappa}^{gfm}
\le q_{i,\xi,\kappa}^{bess} &\le
\overline Q_i^{gfm} z_{i,\xi,\kappa}^{gfm},
\qquad \forall i\in\mathcal{N}^{gfm},
\label{eq:10b-ext-gfm-q}\\
p_{i,\xi,\kappa}^{bess,dis}\tau_{\xi,\kappa} &\le s_i^{ext}\overline E_i^{bess},
\qquad \forall i\in\mathcal{N}^{gfm}.
\label{eq:10b-ext-gfm-energy}
\end{align}
Thus a BESS can form a partition only when it is explicitly activated as a grid-forming source and has sufficient discharge energy to support that partition during the extreme event.

\subsubsection{Extreme-state electrical and load-service constraints}

Conditional on \(\beta_{a,\xi,\kappa}\), the AC/DC power-flow relations are enforced explicitly. For AC branches \(e=(i,j)\in\mathcal{E}^{ac}\), DC branches \(\ell=(i,j)\in\mathcal{E}^{dc}\), and converters \(c\in\mathcal{E}^{vsc}\),
\begin{align}
-M_e^{ac}(1-\beta_{e,\xi,\kappa}) 
\le u_{j,\xi,\kappa}^{ac} - u_{i,\xi,\kappa}^{ac}
+ 2\left(r_e^{ac} f_{e,\xi,\kappa}^{ac} + x_e^{ac} q_{e,\xi,\kappa}^{ac}\right)
\le M_e^{ac}(1-\beta_{e,\xi,\kappa}),
\label{eq:10b-ext-ac-drop}\\
-\overline F_e^{ac}\beta_{e,\xi,\kappa} \le f_{e,\xi,\kappa}^{ac} \le \overline F_e^{ac}\beta_{e,\xi,\kappa},
\label{eq:10b-ext-ac-flow}\\
-\overline Q_e^{ac}\beta_{e,\xi,\kappa} \le q_{e,\xi,\kappa}^{ac} \le \overline Q_e^{ac}\beta_{e,\xi,\kappa},
\label{eq:10b-ext-ac-reactive-flow}\\
\underline u_i^{ac} y_{i,\xi,\kappa}^{en} \le u_{i,\xi,\kappa}^{ac} \le \overline u_i^{ac} y_{i,\xi,\kappa}^{en},
\label{eq:10b-ext-ac-voltage}\\
-M_\ell^{dc}(1-\beta_{\ell,\xi,\kappa})
\le f_{\ell,\xi,\kappa}^{dc} - g_\ell^{dc}(v_{i,\xi,\kappa}^{dc}-v_{j,\xi,\kappa}^{dc})
\le M_\ell^{dc}(1-\beta_{\ell,\xi,\kappa}),
\label{eq:10b-ext-dc-drop}\\
-\overline F_\ell^{dc}\beta_{\ell,\xi,\kappa} \le f_{\ell,\xi,\kappa}^{dc} \le \overline F_\ell^{dc}\beta_{\ell,\xi,\kappa},
\label{eq:10b-ext-dc-flow}\\
\underline v_i^{dc} y_{i,\xi,\kappa}^{en} \le v_{i,\xi,\kappa}^{dc} \le \overline v_i^{dc} y_{i,\xi,\kappa}^{en},
\label{eq:10b-ext-dc-voltage}\\
-\overline S_c^{vsc}\beta_{c,\xi,\kappa} \le p_{c,\xi,\kappa}^{vsc} \le \overline S_c^{vsc}\beta_{c,\xi,\kappa},
\label{eq:10b-ext-vsc-active}\\
-\overline\gamma_c^{vsc}\overline S_c^{vsc}\beta_{c,\xi,\kappa}
\le q_{c,\xi,\kappa}^{vsc} \le
\overline\gamma_c^{vsc}\overline S_c^{vsc}\beta_{c,\xi,\kappa}.
\label{eq:10b-ext-vsc-reactive}
\end{align}
For activated AC source buses \(r\in\mathcal{R}^{ac}\) and activated DC source buses \(m\in\mathcal{R}^{dc}\), reference voltages are imposed by
\begin{align}
-M_r^{src}(1-z_{r,\xi,\kappa}^{src})
\le u_{r,\xi,\kappa}^{ac} - u_r^{ref}
\le M_r^{src}(1-z_{r,\xi,\kappa}^{src}),
\qquad \forall r\in\mathcal{R}^{ac},
\label{eq:10b-ext-ac-root-voltage}\\
-M_m^{src}(1-z_{m,\xi,\kappa}^{src})
\le v_{m,\xi,\kappa}^{dc} - v_m^{ref}
\le M_m^{src}(1-z_{m,\xi,\kappa}^{src}),
\qquad \forall m\in\mathcal{R}^{dc}.
\label{eq:10b-ext-dc-root-voltage}
\end{align}
These constraints ensure that de-energized buses collapse to zero voltage, energized buses remain inside admissible envelopes, and only activated substation roots impose voltage references on their restored partitions. For AC buses,
\begin{align}
\sum_{e\in\delta_i^{-}\cap\mathcal{E}^{ac}} f_{e,\xi,\kappa}^{ac}
- \sum_{e\in\delta_i^{+}\cap\mathcal{E}^{ac}} f_{e,\xi,\kappa}^{ac}
+ p_{i,\xi,\kappa}^{inj,ext}
+ p_{i,\xi,\kappa}^{ls}
= 0,
\label{eq:10b-ext-ac-balance}\\
\sum_{e\in\delta_i^{-}\cap\mathcal{E}^{ac}} q_{e,\xi,\kappa}^{ac}
- \sum_{e\in\delta_i^{+}\cap\mathcal{E}^{ac}} q_{e,\xi,\kappa}^{ac}
+ q_{i,\xi,\kappa}^{inj,ext}
+ \eta_i^{Q,load} p_{i,\xi,\kappa}^{ls}
= 0,
\label{eq:10b-ext-ac-reactive-balance}
\end{align}
and for DC buses,
\begin{align}
\sum_{\ell\in\delta_m^{-}\cap\mathcal{E}^{dc}} f_{\ell,\xi,\kappa}^{dc}
- \sum_{\ell\in\delta_m^{+}\cap\mathcal{E}^{dc}} f_{\ell,\xi,\kappa}^{dc}
+ p_{m,\xi,\kappa}^{inj,ext}
+ p_{m,\xi,\kappa}^{ls}
= 0.
\label{eq:10b-ext-dc-balance}
\end{align}
The extreme-state injection terms are written explicitly as
\begin{align}
p_{i,\xi,\kappa}^{inj,ext} :=
&\mathbf{1}_{\{i\in\mathcal{N}^{pv}\}}p_{i,\xi,\kappa}^{pv}
+ \mathbf{1}_{\{i\in\mathcal{N}^{bess}\setminus\mathcal{N}^{gfm}\}}
\left(p_{i,\xi,\kappa}^{bess,dis}-p_{i,\xi,\kappa}^{bess,ch}\right) \\
&+ \mathbf{1}_{\{i\in\mathcal{N}^{gfm}\}}
\left(p_{i,\xi,\kappa}^{bess,dis}-p_{i,\xi,\kappa}^{bess,ch}\right)
+ \mathbf{1}_{\{i\in\mathcal{R}^{ac}\}} z_{i,\xi,\kappa}^{src}
\left(p_{i,\xi,\kappa}^{sub,+}-p_{i,\xi,\kappa}^{sub,-}\right) \\
&+ \sum_{c:i(c)=i} p_{c,\xi,\kappa}^{vsc}
- P_{i,\xi}^{load,ac},
\label{eq:10b-ext-ac-net-injection}\\
q_{i,\xi,\kappa}^{inj,ext} :=
&\mathbf{1}_{\{i\in\mathcal{N}^{pv}\}}q_{i,\xi,\kappa}^{pv}
+ \mathbf{1}_{\{i\in\mathcal{N}^{bess}\}}q_{i,\xi,\kappa}^{bess}
+ \mathbf{1}_{\{i\in\mathcal{R}^{ac}\}} z_{i,\xi,\kappa}^{src} q_{i,\xi,\kappa}^{sub} \\
&+ \sum_{c:i(c)=i} q_{c,\xi,\kappa}^{vsc}
- Q_{i,\xi}^{load,ac},
\label{eq:10b-ext-ac-reactive-net-injection}\\
p_{m,\xi,\kappa}^{inj,ext} :=
&\mathbf{1}_{\{m\in\mathcal{N}^{pv}\}}p_{m,\xi,\kappa}^{pv}
+ \mathbf{1}_{\{m\in\mathcal{N}^{bess}\}}
\left(p_{m,\xi,\kappa}^{bess,dis}-p_{m,\xi,\kappa}^{bess,ch}\right) \\
&+ \mathbf{1}_{\{m\in\mathcal{R}^{dc}\}} z_{m,\xi,\kappa}^{src}
\left(p_{m,\xi,\kappa}^{sub,+}-p_{m,\xi,\kappa}^{sub,-}\right)
+ \sum_{c:k(c)=m} p_{c,\xi,\kappa}^{vsc,dc}
- P_{m,\xi}^{load,dc}.
\label{eq:10b-ext-dc-net-injection}
\end{align}
Thus only activated substations and activated grid-forming BESS units may act as partition roots in the resilience layer, while non-source DER units remain dispatchable support devices inside an energized partition.

Load service is tied to bus energization by
\begin{align}
0 \le p_{n,\xi,\kappa}^{ls} &\le P_{n,\xi}^{load},
\qquad \forall n\in\mathcal{N},
\label{eq:10b-ext-shed-bound}\\
p_{n,\xi,\kappa}^{ls} &\ge P_{n,\xi}^{load}(1-y_{n,\xi,\kappa}^{en}),
\qquad \forall n\in\mathcal{N}.
\label{eq:10b-ext-deenergized-shed}
\end{align}
Hence a bus that is not energized must be fully shed, while an energized bus may still experience partial curtailment if local generation and transfer capacity are insufficient.

The corresponding extreme-state unsupplied-energy functional is
\begin{align}
H^{ext}_{\xi,\kappa}(x) := \sum_{n\in\mathcal{N}} \upsilon_n^{ext} p_{n,\xi,\kappa}^{ls},
\label{eq:10b-ext-shed-functional}
\end{align}
where \(\upsilon_n^{ext}\) is a load-priority weight that may emphasize critical-service buses during severe events. Constraint \eqref{eq:10b-master-extreme} then limits the weighted resilience loss across the extreme-state library.

This resilience layer is a MILP because it combines binary topology variables, source-activation decisions, and grid-forming mode selection. It extends the peak \(N{-}1\) reliability screen to a restoration-aware planning surrogate for higher-order faults.

\subsection{Discussion of Modeling Choices}

The 10b model is deliberately asymmetric in time resolution:
\begin{enumerate}
  \item carbon accounting and operating cost are driven by annual representative states;
  \item peak-time \(N{-}1\) reliability is screened at the annual stress point, because feeder adequacy violations are typically most severe there;
  \item extreme-state \(N{-}k\) resilience is screened on a compact event library with restoration-aware partitioning;
  \item FMEA weighting and extreme-state severity weighting preserve contingency information without requiring full chronological multi-contingency simulation.
\end{enumerate}

This structure makes the model suitable for distribution-level planning studies where the main design variables are DER siting and sizing, the carbon target is annual, the reliability target is peak adequacy under single failures, and the resilience target is survivability under a limited number of severe higher-order events.

\begin{remark}[Interpretation of the reliability layer]
The peak-only \(N{-}1\) screening layer should still be interpreted as a planning-security approximation rather than a full reliability-assessment replacement. The new extreme-state layer partly closes that gap by introducing restoration-aware \(N{-}k\) events, but it remains a reduced event library rather than a full sequential Monte Carlo simulation. If seasonal restoration behavior is also important, additional event states can be added to \(\Xi\) without changing the model structure.
\end{remark}

\subsection{Illustrative Test System for Case 10b}

To support reproducible numerical experiments, the planning model above is paired with a synthetic but engineering-oriented hybrid AC/DC distribution benchmark. The test system is intended to be large enough to expose the coordination value of AC/DC coupling and DER siting, while remaining small enough to support repeated sensitivity runs.

\subsubsection{Network structure}

The proposed 10b benchmark is a radial hybrid feeder with the following topology.
\begin{align}
&\text{AC buses: } |
\mathcal{N}^{ac}| = 33, \\
&\text{DC buses: } |\mathcal{N}^{dc}| = 9, \\
&\text{AC branches: } |\mathcal{E}^{ac}| = 32, \\
&\text{DC branches: } |\mathcal{E}^{dc}| = 8, \\
&\text{interlinking converters: } |\mathcal{E}^{vsc}| = 3.
\end{align}

The AC side is organized as a modified radial feeder derived from a standard medium-voltage distribution template, so voltage regulation can be represented explicitly through LinDistFlow-type drop equations. The DC side represents converter-fed industrial or campus subzones with strong electronic load penetration. One AC root bus and one DC root bus are connected to the upstream utility through import-limited substations. The three AC/DC converters link the AC backbone to three DC pockets positioned at upstream, middle-feeder, and downstream locations, respectively.

The network is partitioned into three load regions:
\begin{enumerate}
  \item an upstream mixed commercial region with high daytime PV coincidence;
  \item a mid-feeder residential region with evening peak stress;
  \item a downstream critical-service region including hospital, communication, and emergency-support loads.
\end{enumerate}

The critical-service region defines the set \(\mathcal{N}^{crit}\) used in the strict \(N{-}1\) reliability constraints.

\subsubsection{Candidate DER placement}

Candidate PV sites are placed at buses with suitable rooftop or community-solar potential, while candidate BESS sites are placed at electrically strategic buses that can support both congestion relief and contingency restoration. A practical initial candidate set is:
\begin{align}
|\mathcal{N}^{pv}| &= 10,
&|\mathcal{N}^{bess}| &= 6.
\end{align}

The candidate buses should be selected so that at least:
\begin{enumerate}
  \item two PV candidates are near the AC root,
  \item four PV candidates are distributed over middle and downstream AC buses,
  \item two PV candidates are located on DC buses behind converters,
  \item two BESS candidates are placed in the critical-service region,
  \item one BESS candidate is placed near each converter-coupled DC pocket.
\end{enumerate}

This placement pattern allows the planning model to test whether converter capacity, local storage, feeder siting, and grid-forming source placement act as substitutes or complements under carbon, reliability, and resilience pressure.

\subsubsection{Representative operating states}

The annual operating layer uses a compact representative-state library rather than full chronological simulation. A suitable baseline set is:
\begin{align}
\Omega = \{&\text{summer peak},\ \text{summer shoulder},\ \text{winter peak},\ \text{winter shoulder},\\
&\text{spring high-PV},\ \text{autumn high-PV},\ \text{night valley},\ \omega^{pk}\}.
\end{align}

The peak snapshot \(\omega^{pk}\) is the annual system peak and is also the anchor state for the \(N{-}1\) contingency layer. The remaining snapshots are used to approximate annual operating cost and annual carbon exposure. Their weights should be calibrated from feeder smart-meter and PV-availability data or from a synthetic seasonal load/PV library.

\subsubsection{Extreme-state library and \texorpdfstring{$N{-}k$}{N-k} families}

The resilience layer should use a compact but explicit extreme-state library such as
\begin{align}
\Xi = \{&\text{heat-wave import stress},\ \text{storm-driven multi-line outage},\ \text{flooded downstream corridor}\}.
\end{align}
For each \(\xi\in\Xi\), the family \(\mathcal{K}_\xi\) should include geographically or functionally coupled higher-order failures, for example:
\begin{enumerate}
  \item simultaneous outage of two to four adjacent AC feeder sections,
  \item joint outage of a converter and its downstream DC branch pocket,
  \item common-mode loss of one substation and one to two nearby feeders,
  \item storm-induced outage of one corridor plus one large DER block.
\end{enumerate}
For a first reproducible benchmark, it is sufficient to instantiate three event classes with explicit stress assumptions:
\begin{enumerate}
  \item \textbf{heat-wave import stress:} upstream import capability reduced by 20--40\%, annual-peak-like demand retained, and one additional feeder outage added to mimic protection-triggered overload isolation;
  \item \textbf{storm-driven corridor damage:} two to four adjacent AC lines and one converter or DC tie simultaneously unavailable in the same geographical corridor;
  \item \textbf{flooded downstream corridor:} one downstream substation path and all vulnerable downstream laterals inside the flooded zone unavailable, requiring local partition support for critical loads.
\end{enumerate}
The purpose of \(\Xi\) is not exhaustive event enumeration, but a planning-grade set of severe conditions that activate islanding, reactive support scarcity, and multi-source partition restoration.

\subsubsection{Contingency set for peak-time FMEA}

The \(N{-}1\) contingency set \(\mathcal{C}\) should include all single failures whose loss can materially affect peak feasibility:
\begin{align}
\mathcal{C} =
&\mathcal{C}^{ac\_line}
\cup \mathcal{C}^{dc\_line}
\cup \mathcal{C}^{vsc}
\cup \mathcal{C}^{sub}
\cup \mathcal{C}^{der},
\end{align}
where:
\begin{enumerate}
  \item \(\mathcal{C}^{ac\_line}\) contains all primary AC feeder branches;
  \item \(\mathcal{C}^{dc\_line}\) contains all DC feeder links;
  \item \(\mathcal{C}^{vsc}\) contains the three AC/DC converters;
  \item \(\mathcal{C}^{sub}\) contains the AC and DC root-supply outages;
  \item \(\mathcal{C}^{der}\) optionally contains large BESS outages and the largest individual PV block outage.
\end{enumerate}

For the base 10b study, it is sufficient to start with line, converter, and substation contingencies only, then add major DER outages in a secondary sensitivity study.

\subsection{Baseline Parameter Set}

Tables~\ref{tab:10b-parameter-baseline-a}--\ref{tab:10b-parameter-baseline-b} collect a compact baseline parameter set for the first full numerical study. The numbers below are not meant to be universal constants; they provide a coherent starting point for calibration and sensitivity analysis.

\begin{table}[t]
\centering
\caption{Baseline planning and economic parameters for the 10b hybrid AC/DC distribution-planning benchmark.}
\label{tab:10b-parameter-baseline-a}
\begin{tabular}{L{0.34\textwidth}L{0.22\textwidth}L{0.34\textwidth}}
\toprule
Quantity & Baseline value & Interpretation \\
\midrule
Annual discount rate \(r\) & 8\% & annualization basis for DER and converter investment \\
Planning horizon & 1 representative year & annualized planning objective \\
Annual carbon budget \(\overline E^{ann}_{CO_2}\) & 0.80 of base-case emissions & moderate decarbonization target \\
Peak FMEA reliability budget \(\overline E^{pk}_{rel}\) & 0 for critical buses; 0.5--2.0 MWh for all-load sensitivities & strict critical-load security, softer systemwide adequacy options \\
Extreme resilience budget \(\overline E^{ext}_{rel}\) & scenario-weighted 1.0--5.0 MWh equivalent & admissible loss under severe \(N{-}k\) events \\
PV CAPEX & 650--900 \$/kW & installed AC-side equivalent PV cost \\
BESS power CAPEX & 120--180 \$/kW & inverter and power-electronics component \\
BESS energy CAPEX & 180--260 \$/kWh & cell and container energy component \\
Converter CAPEX & 90--150 \$/kW & AC/DC interlinking converter sizing cost \\
Substation expansion CAPEX & 60--120 \$/kW & upstream import-capacity reinforcement \\
PV availability \(\alpha^{pv}_{i,\omega}\) & state-dependent, in \([0,1]\) & representative irradiance profile multiplier \\
BESS SOC bounds \([\underline s_i,\overline s_i]\) & [0.1, 0.9] & normal-state storage operating range \\
Peak pre-contingency SOC \(s_i^{pk}\) & 0.6--0.8 & reliability-prepared storage state at annual peak \\
Extreme-event pre-contingency SOC \(s_i^{ext}\) & 0.7--0.9 & resilience-prepared energy state for grid-forming support \\
Upstream purchase tariff & time-state dependent & used in \eqref{eq:10b-normal-obj} \\
Export tariff & 0.2--0.6 of import tariff & limited compensation for reverse power flow \\
Emission factor \(\kappa_{r,\omega}\) & 0.45--0.70 kg/kWh & upstream electricity carbon intensity \\
Failure frequency \(\lambda_k\) & component-specific & annual contingency occurrence rate from FMEA catalog \\
Outage duration \(\tau_k\) & component-specific & equivalent switching plus repair duration at peak \\
\bottomrule
\end{tabular}
\end{table}

\begin{table}[t]
\centering
\caption{Baseline electrical and resilience-support parameters for the 10b hybrid AC/DC distribution-planning benchmark.}
\label{tab:10b-parameter-baseline-b}
\begin{tabular}{L{0.34\textwidth}L{0.22\textwidth}L{0.34\textwidth}}
\toprule
Quantity & Baseline value & Interpretation \\
\midrule
PV reactive capability factor \(\overline\gamma_i^{pv}\) & 0.3--0.5 p.u. & inverter VAR capability relative to available active power \\
BESS reactive capability factor \(\overline\gamma_i^{bess}\) & 0.4--0.6 p.u. & inverter VAR capability relative to BESS power rating \\
VSC reactive capability factor \(\overline\gamma_c^{vsc}\) & 0.3--0.6 p.u. & AC-side converter VAR capability factor \\
AC line resistance \(r_e^{ac}\) & feeder-specific & active-power contribution to LinDistFlow voltage drop \\
AC line reactance \(x_e^{ac}\) & feeder-specific & reactive-power contribution to LinDistFlow voltage drop \\
AC active flow limits & feeder-specific & branch active-power bounds under AC P/Q LinDistFlow \\
AC reactive flow limits & feeder-specific & branch reactive-power bounds under AC P/Q LinDistFlow \\
AC voltage bounds \([\underline u_i^{ac},\overline u_i^{ac}]\) & feeder-specific, typically 0.95\(^2\)--1.05\(^2\) p.u. & squared-voltage operating envelope for AC buses \\
AC root voltage \(u_r^{ref}\) & 1.0\(^2\) p.u. & feeder-head LinDistFlow reference \\
AC root reactive capability \(\overline Q_r^{sub}\) & feeder-specific & upstream reactive support range at each AC substation \\
AC load reactive ratio \(\eta_i^{Q,load}\) & bus-specific or region-specific & converts active load service and curtailment into reactive demand \\
DC line thermal limits & feeder-specific & branch active-power limits under linear DC flow \\
Contingency decoupling constants \(M_e^{ac},M_\ell^{dc}\) & sufficiently large, topology-specific & deactivate constitutive equations for outaged branches \\
Grid-forming BESS active capability \(\overline P_i^{gfm}\) & subset-specific & active-power headroom available in islanded source mode \\
Grid-forming BESS reactive capability \(\overline Q_i^{gfm}\) & subset-specific & voltage support capability when BESS forms a partition \\
Extreme-state weights \(\varpi_\xi\) & state-specific & planning importance or occurrence weight of each severe event \\
Extreme-event durations \(\tau_{\xi,\kappa}\) & contingency-specific & effective restoration duration for each \(N{-}k\) event \\
\bottomrule
\end{tabular}
\end{table}

The main purpose of Tables~\ref{tab:10b-parameter-baseline-a}--\ref{tab:10b-parameter-baseline-b} is to freeze a reporting basis. Once numerical runs begin, exact values should be mapped to the component fields already documented in Section~\ref{sec:distribution-reliability-fmea} and the I/O contract chapters.

\subsection{Experimental Design}

The numerical study for 10b should answer four questions simultaneously:
\begin{enumerate}
  \item how the hybrid feeder uses DER siting and converter capacity to reduce annual carbon exposure;
  \item how much additional capacity is required to satisfy peak-time \(N{-}1\) reliability targets;
  \item how much DER-assisted voltage regulation reduces reinforcement and load curtailment under stressed conditions;
  \item whether AC/DC coupling and grid-forming BESS change survivability under extreme \(N{-}k\) events relative to a pure AC planning baseline.
\end{enumerate}

To answer those questions cleanly, the experiments should be organized into a small number of structured cases.

\subsubsection{Evaluation metrics}

Every case should report the same output set:
\begin{align}
&\text{installed PV capacity by bus and total}, \\
&\text{installed BESS power and energy by bus and total}, \\
&\text{installed converter and substation capacities}, \\
&\text{planning objective and its investment / operating decomposition}, \\
&\text{annual carbon emission and carbon-budget slack}, \\
&\text{peak FMEA-weighted EENS and contingency-wise shed-load diagnostics}, \\
&\text{extreme-state weighted ENS and partition-count diagnostics}, \\
&\text{critical-load service status under every } N{-}1 \text{ contingency}, \\
&\text{critical-load service status under every extreme } N{-}k \text{ event}, \\
&\text{root import dependence of the AC and DC subnetworks}, \\
&\text{solver runtime and model size}. 
\end{align}

In addition, the following structural diagnostics are recommended:
\begin{align}
&\text{utilization of each converter at } \omega^{pk}, \\
&\text{fraction of peak contingency support supplied by local BESS}, \\
&\text{change in downstream branch loading after DER installation}, \\
&\text{number and composition of energized tree partitions under each extreme event}, \\
&\text{reactive-support contribution shares of PV, BESS, VSC, and substations}. 
\end{align}

\subsubsection{Case 1: Base AC/DC planning without carbon or reliability pressure}

The first case is the economic reference point:
\begin{align}
&\text{annual operating states enabled}, \\
&\zeta^{CO_2} \text{ inactive or budget very loose}, \\
&\zeta^{rel} \text{ inactive or reliability budget very loose}, \\
&\zeta^{ext} \text{ inactive or resilience budget very loose}, \\
&\text{lossless converter model enabled}. 
\end{align}

This case identifies the purely economic DER pattern selected by the hybrid feeder without policy or security pressure. It serves as the denominator for all later comparisons.

\subsubsection{Case 2: Carbon-constrained planning}

The second case activates the annual carbon constraint \eqref{eq:10b-master-carbon} while keeping the reliability layer soft or inactive. Its purpose is to quantify how much PV, BESS, and converter reinforcement are induced by decarbonization alone.

Recommended sensitivities are:
\begin{align}
\overline E^{ann}_{CO_2} \in \{1.0,\ 0.9,\ 0.8,\ 0.7\}\times E^{ann}_{CO_2,\mathrm{base}}.
\end{align}

The main outputs of interest are marginal carbon abatement cost, relocation of DER from AC to DC buses, and the extent to which converter expansion substitutes for local storage.

\subsubsection{Case 3: Peak-time \texorpdfstring{$N{-}1$}{N-1} reliability-constrained planning}

The third case activates the peak FMEA constraint \eqref{eq:10b-fmea-peak-eens} and the strict critical-load condition \eqref{eq:10b-critical-n-minus-1}, while keeping the carbon budget loose. This isolates the planning cost of reliability.

Two variants are especially informative:
\begin{enumerate}
  \item \textbf{critical-load security only:} \(p_{n,k}^{ls}=0\) only on \(\mathcal{N}^{crit}\);
  \item \textbf{systemwide adequacy:} \(\overline E^{pk}_{rel}\) tightened for all buses.
\end{enumerate}

This case reveals whether the model prefers local BESS in the critical-service region, stronger upstream import capacity, additional converter reinforcement, or more flexible inverter VAR support to survive single contingencies at peak.

\subsubsection{Case 4: Joint carbon- and reliability-constrained planning}

The fourth case activates \eqref{eq:10b-master-carbon} and \eqref{eq:10b-master-reliability} together with DER-assisted voltage regulation. This isolates the value of coordinated PV/BESS/VSC VAR support before the resilience layer is added.

The most useful outputs are not only total cost and total capacity, but also the interaction effects:
\begin{enumerate}
  \item whether carbon-constrained PV siting improves or worsens peak \(N{-}1\) feasibility;
  \item whether reliability-driven BESS placement also reduces carbon exposure by smoothing root imports;
  \item whether DER-assisted voltage regulation reduces the reinforcement needed to sustain the carbon-reliability combination. 
\end{enumerate}

\subsubsection{Case 5: Extreme-state \texorpdfstring{$N{-}k$}{N-k} resilience planning}

The fifth case activates the full model, including the resilience constraint \eqref{eq:10b-master-extreme}, source-partition binaries, and grid-forming BESS constraints. This is the main resilience case.

The most informative outputs are:
\begin{enumerate}
  \item how often the model chooses multiple energized partitions rather than a single surviving backbone,
  \item whether grid-forming BESS is concentrated in the critical-service region or distributed across the feeder,
  \item how much extreme-event ENS is reduced by reactive support and converter flexibility,
  \item whether extreme-state preparedness changes the carbon and peak \(N{-}1\) design selected by the first stage.
\end{enumerate}

\subsubsection{Case 6: Architecture and device-value sensitivities}

To isolate the value of hybridization itself, the integrated resilience-aware design should be compared against the following structural ablations:
\begin{align}
&\text{pure AC benchmark: } \mathcal{N}^{dc}=\varnothing, \\
&\text{fixed-converter benchmark: } \overline S_c^{vsc} \text{ not expandable}, \\
&\text{no-BESS benchmark: } \overline P_i^{bess}=\overline E_i^{bess}=0, \\
&\text{no-DER benchmark: only substation reinforcement allowed}. 
\end{align}

These ablations show whether the main benefit arises from topology, from storage, from inverter VAR support, from grid-forming capability, or from the AC/DC coupling device itself.

\subsection{Recommended Reporting Tables and Figures}

For a first complete 10b numerical study, the following outputs should be prepared.

\begin{enumerate}
  \item A system-data table listing buses, candidate DER sites, critical-load buses, converters, and contingency counts.
  \item Baseline-parameter tables, starting from Tables~\ref{tab:10b-parameter-baseline-a}--\ref{tab:10b-parameter-baseline-b}.
  \item A case-summary table comparing Cases 1--6 in objective value, carbon emission, peak FMEA-EENS, extreme-state ENS, and installed capacities.
  \item A bus-level DER-siting figure showing where PV and BESS are built under the joint case.
  \item A contingency heatmap showing peak shed load by contingency and by load region.
  \item An extreme-event partition figure showing energized islands and their source buses under one representative \(N{-}k\) event.
  \item A carbon--reliability--resilience tradeoff plot obtained by varying \(\overline E^{ann}_{CO_2}\), \(\overline E^{pk}_{rel}\), and \(\overline E^{ext}_{rel}\).
\end{enumerate}

\subsection{Expected Experimental Findings}

Although the exact numerical values must be determined computationally, the 10b model is designed to test five hypotheses.

\begin{enumerate}
  \item Hybrid AC/DC coupling should reduce the cost of meeting peak \(N{-}1\) reliability relative to a pure AC feeder because mobile power support across subnetworks is enabled through the interlinking converters.
  \item Carbon tightening should increase PV adoption first, then BESS adoption, and finally converter/substation reinforcement once local renewable hosting limits become active.
  \item Reliability tightening should shift BESS placement toward downstream critical-service buses and increase the option value of converter capacity even when converters are modeled as lossless.
  \item DER-assisted voltage regulation should reduce both peak-time curtailment and the need for substation VAR expansion compared with a substation-only reactive support policy.
  \item The extreme-state resilience layer should induce additional BESS energy reservation and more spatially distributed inverter placement, because multiple tree-structured partitions become valuable under severe \(N{-}k\) events.
\end{enumerate}

These hypotheses provide a clear roadmap for the first computational campaign and convert the 10b section from a pure formulation note into an executable numerical-study plan.

\subsection{Result-Display Skeleton}

To reduce setup overhead in the first numerical campaign, the following reporting templates are prepared directly in the notebook workspace.

\subsubsection{Case-summary and parameter-table templates}

The case-summary table template is stored in
\texttt{visualization/hybrid\_distribution\_case10b/tables/table\_10b\_case\_summary.tex}, and the compact parameter-summary template is stored in
\texttt{visualization/hybrid\_distribution\_case10b/tables/table\_10b\_parameter\_summary.tex}. Their CSV companions are stored in the matching \texttt{data/} directory so that result-generation scripts can populate the same reporting layout automatically.

\input{visualization/hybrid_distribution_case10b/tables/table_10b_case_summary.tex}

\input{visualization/hybrid_distribution_case10b/tables/table_10b_parameter_summary.tex}

\subsubsection{Figure placeholders}

Three figure placeholders are also prepared for the first paper draft: feeder topology, carbon--reliability tradeoff, and DER siting under the joint case.

\input{visualization/hybrid_distribution_case10b/figures/fig_10b_topology_placeholder.tex}

\input{visualization/hybrid_distribution_case10b/figures/fig_10b_case_tradeoff_placeholder.tex}

\input{visualization/hybrid_distribution_case10b/figures/fig_10b_der_siting_placeholder.tex}